Tentukan semua fungsi $f:N\rightarrow N$ yang memenuhi
$$f(mn)+f(m+n)=f(m)f(n)+1$$
untuk semua $m, n\in N$.
\begin{remark*}[Motivasi dan Walkthrough]

\begin{itemize}
    \item $2+2 = 2\times2$.
    \item Karena fungsi bulat, maka bisa coba induksi.
    \item Coba substitusi $m=1$ atau $n=1$ (standar lah ya karena domainnya bulat positif).
    \item Ada "+1" di akhir harusnya curiga bisa "dihilangkan".
\end{itemize}
\end{remark*}
\begin{solusi}

Misalkan $P(m,n)$ adalah asersi $f(mn)+f(m+n)=f(m)f(n)+1$.

Kita punya
\begin{align*}
    P(m,1): f(m)+f(m+1)=f(m)f(1)+1 \implies f(m+1)=(f(1)-1)f(m)+1.
\end{align*}
Hal yang cukup menarik untuk diulik adalah nilai $(f(1)-1)f(m)$ di RHS, apakah akan nol atau tidak. Oleh karena itu bisa dibagi kasus seperti berikut.\\

\textbf{Kasus 1: $f(1)=1$}\\
Maka $f(m+1)=1$ untuk semua $m\in N$. Jadi $f(m)=1$ untuk semua $m\in N$. Cek ke persamaan, masih memenuhi.\\


\textbf{Kasus 2: $f(1)\neq 1$ atau $f(1)>1$ karena $f(1)\in N$}\\
Maka dapat dimisalkan $f(1)-1=k$ untuk suatu bilangan bulat positif konstan $k$. Maka $f(m+1)=kf(m)+1$. 
Oleh karena itu, dengan induksi mudah didapatkan bahwa 
\begin{align*}
    f(m)=k^{m}+k^{m-1}+\cdots+k+1
\end{align*}
untuk semua $m\in N$.

Namun, di sisi lain kita punya
\begin{align*}
    P(2,2): f(4)+f(4)&=f(2)f(2)+1 \\
    \implies 2f(4)&=f(2)^2+1\\
    \implies 2(k^4+k^3+k^2+k+1)&=(k^2+k+1)^2+1\\
    \implies 2k^4+2k^3+2k^2+2k+2&=k^4+2k^3+3k^2+2k+2\\
    \implies k^4-k^2&=0\\
    \implies k^2(k^2-1)&=0\\
    \implies k&=1 \text{ Karena $k\in N$}.
\end{align*}
Berarti didapat $f(m)=1+1+\cdots+1= m+1$ untuk semua $m\in N$. Cek ke persamaan, fungsi ini juga memenuhi.

Dari kedua kasus di atas, maka diperoleh $f(m)=1$ atau $f(m)=m+1$ untuk semua $m\in N$.
\end{solusi}